% Template for PLoS
% Version 3.5 March 2018
%
% % % % % % % % % % % % % % % % % % % % % %
%
% -- IMPORTANT NOTE
%
% This template contains comments intended
% to minimize problems and delays during our production
% process. Please follow the template instructions
% whenever possible.
%
% % % % % % % % % % % % % % % % % % % % % % %
%
% Once your paper is accepted for publication,
% PLEASE REMOVE ALL TRACKED CHANGES in this file
% and leave only the final text of your manuscript.
% PLOS recommends the use of latexdiff to track changes during review, as this will help to maintain a clean tex file.
% Visit https://www.ctan.org/pkg/latexdiff?lang=en for info or contact us at latex@plos.org.
%
%
% There are no restrictions on package use within the LaTeX files except that
% no packages listed in the template may be deleted.
%
% Please do not include colors or graphics in the text.
%
% The manuscript LaTeX source should be contained within a single file (do not use \input, \externaldocument, or similar commands).
%
% % % % % % % % % % % % % % % % % % % % % % %
%
% -- FIGURES AND TABLES
%
% Please include tables/figure captions directly after the paragraph where they are first cited in the text.
%
% DO NOT INCLUDE GRAPHICS IN YOUR MANUSCRIPT
% - Figures should be uploaded separately from your manuscript file.
% - Figures generated using LaTeX should be extracted and removed from the PDF before submission.
% - Figures containing multiple panels/subfigures must be combined into one image file before submission.
% For figure citations, please use "Fig" instead of "Figure".
% See http://journals.plos.org/plosone/s/figures for PLOS figure guidelines.
%
% Tables should be cell-based and may not contain:
% - spacing/line breaks within cells to alter layout or alignment
% - do not nest tabular environments (no tabular environments within tabular environments)
% - no graphics or colored text (cell background color/shading OK)
% See http://journals.plos.org/plosone/s/tables for table guidelines.
%
% For tables that exceed the width of the text column, use the adjustwidth environment as illustrated in the example table in text below.
%
% % % % % % % % % % % % % % % % % % % % % % % %
%
% -- EQUATIONS, MATH SYMBOLS, SUBSCRIPTS, AND SUPERSCRIPTS
%
% IMPORTANT
% Below are a few tips to help format your equations and other special characters according to our specifications. For more tips to help reduce the possibility of formatting errors during conversion, please see our LaTeX guidelines at http://journals.plos.org/plosone/s/latex
%
% For inline equations, please be sure to include all portions of an equation in the math environment.
%
% Do not include text that is not math in the math environment.
%
% Please add line breaks to long display equations when possible in order to fit size of the column.
%
% For inline equations, please do not include punctuation (commas, etc) within the math environment unless this is part of the equation.
%
% When adding superscript or subscripts outside of brackets/braces, please group using {}.
%
% Do not use \cal for caligraphic font.  Instead, use \mathcal{}
%
% % % % % % % % % % % % % % % % % % % % % % % %
%
% Please contact latex@plos.org with any questions.
%
% % % % % % % % % % % % % % % % % % % % % % % %

\documentclass[10pt,letterpaper]{article}
\usepackage[top=0.85in,left=2.75in,footskip=0.75in]{geometry}

% amsmath and amssymb packages, useful for mathematical formulas and symbols
\usepackage{amsmath,amssymb}

% Use adjustwidth environment to exceed column width (see example table in text)
\usepackage{changepage}

% Use Unicode characters when possible
\usepackage[utf8x]{inputenc}

% textcomp package and marvosym package for additional characters
\usepackage{textcomp,marvosym}

% cite package, to clean up citations in the main text. Do not remove.
% \usepackage{cite}

% Use nameref to cite supporting information files (see Supporting Information section for more info)
\usepackage{nameref,hyperref}

% line numbers
\usepackage[right]{lineno}

% ligatures disabled
\usepackage{microtype}
\DisableLigatures[f]{encoding = *, family = * }

% color can be used to apply background shading to table cells only
\usepackage[table]{xcolor}

% array package and thick rules for tables
\usepackage{array}

% create "+" rule type for thick vertical lines
\newcolumntype{+}{!{\vrule width 2pt}}

% create \thickcline for thick horizontal lines of variable length
\newlength\savedwidth
\newcommand\thickcline[1]{%
  \noalign{\global\savedwidth\arrayrulewidth\global\arrayrulewidth 2pt}%
  \cline{#1}%
  \noalign{\vskip\arrayrulewidth}%
  \noalign{\global\arrayrulewidth\savedwidth}%
}

% \thickhline command for thick horizontal lines that span the table
\newcommand\thickhline{\noalign{\global\savedwidth\arrayrulewidth\global\arrayrulewidth 2pt}%
\hline
\noalign{\global\arrayrulewidth\savedwidth}}


% Remove comment for double spacing
%\usepackage{setspace}
%\doublespacing

% Text layout
\raggedright
\setlength{\parindent}{0.5cm}
\textwidth 5.25in
\textheight 8.75in

% Bold the 'Figure #' in the caption and separate it from the title/caption with a period
% Captions will be left justified
\usepackage[aboveskip=1pt,labelfont=bf,labelsep=period,justification=raggedright,singlelinecheck=off]{caption}
\renewcommand{\figurename}{Fig}

% Use the PLoS provided BiBTeX style
% \bibliographystyle{plos2015}

% Remove brackets from numbering in List of References
\makeatletter
\renewcommand{\@biblabel}[1]{\quad#1.}
\makeatother



% Header and Footer with logo
\usepackage{lastpage,fancyhdr,graphicx}
\usepackage{epstopdf}
%\pagestyle{myheadings}
\pagestyle{fancy}
\fancyhf{}
%\setlength{\headheight}{27.023pt}
%\lhead{\includegraphics[width=2.0in]{PLOS-submission.eps}}
\rfoot{\thepage/\pageref{LastPage}}
\renewcommand{\headrulewidth}{0pt}
\renewcommand{\footrule}{\hrule height 2pt \vspace{2mm}}
\fancyheadoffset[L]{2.25in}
\fancyfootoffset[L]{2.25in}
\lfoot{\today}

%% Include all macros below

\newcommand{\lorem}{{\bf LOREM}}
\newcommand{\ipsum}{{\bf IPSUM}}






\usepackage{forarray}
\usepackage{xstring}
\newcommand{\getIndex}[2]{
  \ForEach{,}{\IfEq{#1}{\thislevelitem}{\number\thislevelcount\ExitForEach}{}}{#2}
}

\setcounter{secnumdepth}{0}

\newcommand{\getAff}[1]{
  \getIndex{#1}{}
}

\providecommand{\tightlist}{%
  \setlength{\itemsep}{0pt}\setlength{\parskip}{0pt}}

\begin{document}
\vspace*{0.2in}

% Title must be 250 characters or less.
\begin{flushleft}
{\Large
\textbf\newline{Prediction of PIK3CA mutation with gene expression} % Please use "sentence case" for title and headings (capitalize only the first word in a title (or heading), the first word in a subtitle (or subheading), and any proper nouns).
}
\newline
% Insert author names, affiliations and corresponding author email (do not include titles, positions, or degrees).
\\
Jun Kang\textsuperscript{\getAff{Department of Hospital Pathology, Seoul St.~Mary's Hospital, College of
Medicine, The Catholic University of Korea, Seoul, South Korea}},
Ahwon Lee\textsuperscript{\getAff{Department of Hospital Pathology, Seoul St.~Mary's Hospital, College of
Medicine, The Catholic University of Korea, Seoul, South Korea}},
Youn Soo Lee\textsuperscript{\getAff{Department of Hospital Pathology, Seoul St.~Mary's Hospital, College of
Medicine, The Catholic University of Korea, Seoul, South Korea}}\textsuperscript{*}\\
\bigskip
\bigskip
* Corresponding author: lys9908@catholic.ac.kr\\
\end{flushleft}
% Please keep the abstract below 300 words
\section*{Abstract}
Breast cancer with PIK3CA mutation has been approved to use PIK3CA
inhibitor in hormone receptor-positive HER2 negative subtype. We apply a
supervised elastic net penalized logistic regression model in prediction
PIK3CA mutation from gene expression data. We used the TCGA pan-cancer
dataset. Penalized logistic regression was applied to prediction
modeling. 10-fold cross-validation with targe variable stratification
was done over the hyperparameter grid. 10845 cases were available for
both PIK3CA mutation and mRNA expression data. In 10-fold
cross-validation, the model with \(\lambda\) = 0.01 and \(\alpha\) = 1.0
(Ridge regression) showed best performance in terms of AUROC. The final
model was trained with the selected hyperparameters with the entire
trainset. The trainset AUROC was 0.93 and the testset AUROC was 0.84.
The AUPR of trainset was 0.66 and the testset AUPR was 0.39. Both IGFR1R
and PTEN are key players in the tyrosine kinase pathway. The cancer type
was an important predictor. Our study suggests that the prediction of
genomic alteration with gene expression data is possible with good
performance.

% Please keep the Author Summary between 150 and 200 words
% Use first person. PLOS ONE authors please skip this step.
% Author Summary not valid for PLOS ONE submissions.

\linenumbers

% Use "Eq" instead of "Equation" for equation citations.
\hypertarget{introduction}{%
\section{Introduction}\label{introduction}}

Targeted therapy becomes a standard treatment in many cancer patients.
Much targeted therapy requires a test for a specific cancer genomic
alteration, to treat the patients. Many direct tests for the genomic
alteration have been developed and prooved for their clinical utility to
find which patients will be responding to the targeted therapy.

It has been actively researched that the machine learnig approach can be
applied to detect genomic alterations. Machine learning algorithm can
build a prediction model from a large number of predictors such as
radiolomic features {[}1{]}, pathology image {[}2{]} or gene expression
data {[}3{]}. Because the most direct genomic test is more specific and
sensitive than predictive models, the machine learning approach might
have a limited role in clinical practice. However, the machine learning
prediction can be a second-best when the direct test fails.

Prediction RAS pathway activation with gene expression data was done in
the previous study. {[}5{]} They trained pan-cancer The Cancer Genome
Atlas (TCGA) data with a supervised elastic net penalized logistic
regression classifier with stochastic gradient descent. The performance
of their model was 84\% with an area under the receiver operating
characteristic (AUROC) curve and 63\% with an area under the
precision-recall (AUPR) curve. The authors suggested their approach can
be applied to other genomic alterations.

PIK3CA encodes the p110\(\alpha\) catalytic subunit of
phosphatidylinositol 3′-kinase (PI3K). PI3K is a protein kinase that
phosphorylates phosphatidylinositol 4,5-biphosphate
(PIP\textsubscript{2}) to make phosphatidylinositol 3,4,5-triphosphate
(PIP\textsubscript{3}). Phosphatase and tensin homolog (PTEN) changes
PIP\textsubscript{2} to PIP\textsubscript{3} in contrast, PI3K.
PIP\textsubscript{3} is a second messenger to activate protein kinase B
(AKT) which is a serine/threonine-specific protein kinase. AKT inhibits
apoptosis and promotes cell proliferation. {[}6{]}

Breast cancer with PIK3CA mutation has been approved to use PIK3CA
inhibitor in hormone receptor-positive HER2 negative subtype. {[}7{]}
The PIK3CA mutation is the second most driver mutation after TP53. The
PIK3CA mutation is most frequently founded in endometrial carcinoma
(45\%), and followed by breast invasive carcinoma (24\%), cervical
squamous cell carcinoma, and endocervical adenocarcinoma (20\%) and
colon adenocarcinoma (16\%).

We apply a supervised elastic net penalized logistic regression model in
prediction PIK3CA mutation. The purpose of this study is to know this
prediction model approach can be applied not only to RAS pathway
activation but also PIK3CA mutation across many cancer types.

\hypertarget{materials-and-methods}{%
\section{Materials and Methods}\label{materials-and-methods}}

\hypertarget{dataset}{%
\subsection{Dataset}\label{dataset}}

We used the TCGA pan-cancer dataset. TCGA is a cancer genomic consortium
that archives data of exome sequencing, gene expression, DNA
methylation, protein expression, and clinical data of more than 10000
cancer samples across 33 common cancer types. The gene expression TCGA
pan-cancer dataset is batch-corrected with normalization. The TCGA
dataset is publically available. PIK3CA mutation data was got using
cgdsr rpackage.{[}8{]} Gene expression data was downloaded from the
National Cancer Institute (NCI)'s Genomic Data Commons (GDC) website
that archives data used for The Pan-Cancer Atlas initiative.
(https://gdc.cancer.gov/about-data/publications/pancanatlas)

The target variable was the PIK3CA mutation status. The status of PIK3CA
was considered as positive when the case has following PIK3CA variants
which are the target variables of the therascreen PIK3CA RGQ PCR Kit;
C420R, E542K, E545A, E545D, E545G, E545K, Q546E, Q546R, H1047L, H1047R,
H1047Y. The therascreen PIK3CA RGQ PCR Kit was approved as a companion
diagnostics to treat with PIK3CA inhibitor by the U.S. Food and Drug
Administration. We split the three-quarters of the dataset for the
trainset and one quarter for the test set.

\hypertarget{modeling-process}{%
\subsection{Modeling process}\label{modeling-process}}

To narrow down potential predictors, genes with a large the median
absolute deviation (more than third-quartiles) were selected. 33 cancer
type dummy variables were included in predictor variables. Yeo-Johnson
transformation was done to correct skewness. Centering and scaling were
done. All preprocessing was done using recipe r package. {[}9{]}
Penalized logistic regression was applied to prediction modeling.
10-fold cross-validation with targe variable stratification was done
over the hyperparameter grid: \(\lambda\) \{\(10^{-5}\),
\(10^{-4}\),\(10^{-3}\),\(10^{-2}\),\(10^{-1}\), \(10^{0}\)\},
\(\alpha\) \{0.0, 0.25, 0.5, 0.75\}. Lambda is penalty scaling parameter
and alpha is mixing parameter of penalty function
(\((1-\alpha)/2 \lVert\beta\rVert_2^2+\alpha\lVert \beta \rVert_1\)).
{[}10{]}

\hypertarget{accessing-model-performance}{%
\subsection{Accessing model
performance}\label{accessing-model-performance}}

Model performance was evaluated with the area under the receiver
operating characteristic (ROC) and area under the ROC (AUROC) and
precision-recall (PR) and area under the curve (AUPR). The AUPR is more
informative than AUROC on imbalanced datasets. {[}11{]} The modeling
process and accessing model performance were done with tidymodels
rpackage. {[}12{]}

\hypertarget{results}{%
\section{Results}\label{results}}

\hypertarget{dataset-summary}{%
\subsection{Dataset summary}\label{dataset-summary}}

10845 cases were available for both PIK3CA mutation and mRNA expression
data. 5128 out of 20502 genes were included in the modeling process
after filtering with median absolute deviation as described at the
modeling process method. The prevalence rate of PIK3CA was 0.11 in all
cases. The PIK3CA prevalence rate of each cancer type varied. The median
prevalence rate of PIK3CA of each cancer type was 0.03 (range 0-0.33)
(Figure 1).

\hypertarget{selecting-model-and-performance-estimation}{%
\subsection{Selecting model and performance
estimation}\label{selecting-model-and-performance-estimation}}

In 10-fold cross-validation, the model with \(\lambda\) = 0.01 and
\(\alpha\) = 1.0 (Ridge regression) showed best performance in terms of
AUROC. The final model was trained with the selected hyperparameters
with the entire trainset. The trainset AUROC was 0.93 and the testset
AUROC was 0.84. The AUPR of trainset was 0.66 and the testset AUPR was
0.39. (Figure 2A)

\hypertarget{performance-of-each-cancer-type}{%
\subsection{Performance of each cancer
type}\label{performance-of-each-cancer-type}}

Because the prevalence of PIK3CA mutation varies across the cancer type,
the performance of each cancer type was investigated. The AUROC and AUPR
were positively correlated between the train set and test set in cancer
type subanalysis. (Figure 2B) The AUPR was high in cancer type with high
PIK3CA mutation rate such as colon, breast, uterus cancer types. The
AUROC did not correlate with the PIK3CA mutation rate of each cancer
type. (Figure 2C)

\hypertarget{important-predictors}{%
\subsection{Important predictors}\label{important-predictors}}

Figure 3 shows the top 30 important predictors. The coefficient is the
parameter of the predictor which represents the effect of the predictor
on prediction. IGF1R mRNA expression was the strongest negative
predictor and PTEN was the strongest positive predictor. Both IGFR1R and
PTEN are key players in the tyrosine kinase pathway. The cancer type was
an important predictor. Some cancer types including uterine
carcinosarcoma (UCS), bladder urothelial carcinoma (BLCA), pancreatic
adenocarcinoma (PAAD), lymphoid neoplasm diffuse large B-cell lymphoma
(DLBC) are the strongest predictors.

\hypertarget{discussion}{%
\section{Discussion}\label{discussion}}

Our model showed good performance to predict the PIK3CA mutation of
various cancer types. This result suggests that the supervised elastic
net penalized logistic regression model can be applied not only to the
RAS activation pathway but also other genomic alterations. Both the RAS
activation pathway and PIK3CA mutation are important and common cancer
genomic alterations. They have a significant effect on gene expression
in cancer cells. It can not be generalized that the supervised elastic
net penalized logistic regression model can be applied to other genomic
alteration. It might be a challenging prediction of genomic alterations
that are infrequent or have a weak effect on gene expression.

Prediction modeling from the TCGA pan-cancer dataset has limitations
regarding data preprocessing. Methods for cancer gene expression have
been developed for research. The gene expression data is processed
between-sample normalization to remove the batch effect. If the model
has been trained from the between-sample normalization, a new sample can
not be exactly preprocessed with normalization which was done on
trainset. A model based on gene expression TCGA pan-cancer dataset has a
limitation on data preprocessing. It is necessary for developing a
preprocessing method that is independent with a dataset to apply the
gene expression data to the prediction model.

Our PIK3CA prediction model performed similar with the RAS activation
prediction model of the previous study in terms of AUROC (0.84). However
AUPR of our model for PIK3CA is lower than model for RAS activation
(0.39 vs 0.63). The reason of lower AUPR of our model can be explained
by low prevalence rate of PIK3CA mutation and imbalanced dataset. The
model for RAS activation trained with cancer types with more than 0.05
prevalence of RAS activation to avoid imbalance classification problem.
We included all cancer types into modeling process. The lower prevalance
rate of target variable mean our dataset has a lower baseline of AUPR.
In the subanalysis for performance of each cancer types, the cancer
types with higher PIK3CA mutation rate showe better AUPR.

Our model includes cancer type predictor and they are stronger than gene
expression data. The varying prevalence of PIC3CA mutation across cancer
type might reason for the strong cancer type predictor. If the cancer
type was wrong or can not be determined, our model performance can be
poor.

Some significant gene expression predictors were closely related to PTEN
and the PI3K pathway. PTEN and IGFR1R are the strongest gene expression
predictor which have negative and positive predictive power. IGF1R is a
tyrosine kinase receptor that activates PI3K. {[}13{]} Insulin receptor
substrate-2 (IRS2) is the adaptor protein of IGF1R. {[}14{]} PTEN is an
important regulator of PIP\textsubscript{3} by dephosphorylating
PIP\textsubscript{3} in contrast PI3K.{[}6{]}

There is some study trying to predict genomic alterations from gene
expression data. A study performed PIK3CA mutation prediction by a
gene-expression signature which is a sum of the average of the
logarithmic gene expression. The model showed good performance AUROC
0.71 in an independent test set. {[}4,15{]} Another study predicted copy
number alterations with gene expression using a multinomial logistic
regression model with least absolute shrinkage and selection operator
(LASSO). The prediction of the 1p/19q codel was very good with an AUROC
of 0.997. The gene-level prediction was good with an AUROC 0.75.
{[}16{]} A Hidden Markov Model Approach for Prediction of Genomic
Alterations from Gene Expression Profiling.{[}17{]} A logistic
regression model was used for MYCN gene amplification in neuroblastoma.
{[}18{]}

Our study suggests that the prediction of genomic alteration with gene
expression data is possible with good performance. However, improvement
of performance is required for applying to clinical test and the
standardization of generation processing of gene expression data is also
needed.

\hypertarget{figure-legends}{%
\section{Figure legends}\label{figure-legends}}

\begin{itemize}
\item
  Figure 1. prevalence rate of PIK3CA across cancer types The
  abbreviations of cancer types are explained in the S1 Appendix.
\item
  Figure 2. Summary of modeling results
\end{itemize}

\begin{enumerate}
\def\labelenumi{(\Alph{enumi})}
\tightlist
\item
  Left: receiver operating characteristic (ROC) curve right:
  precision-recall (PR) curve of trainset and test set. The horizontal
  green line is the PIK3CA mutation rate (0.11) (B) Correlation between
  trainset and test set of the area under the receiver operating
  characteristic curve (AUROC) and area under the precision-recall curve
  (AUPR) among cancer type. The gray band is 95\% confidence interval.
  The abbreviations are explained in the S1 Appendix. (C) Correlation
  between the PIK3CA mutation rate of area under the receiver operating
  characteristic curve (AUROC) and area under the precision-recall curve
  (AUPR).
\end{enumerate}

\begin{itemize}
\tightlist
\item
  Figure 3. Coefficients of model
\end{itemize}

\begin{enumerate}
\def\labelenumi{(\Alph{enumi})}
\tightlist
\item
  Top 30 high coefficients of mRNA. (B) Coefficients of cancer types.
  The abbreviations of cancer types are explained in the S1 Appendix.
\end{enumerate}

\hypertarget{supporting-information}{%
\section{Supporting information}\label{supporting-information}}

\begin{itemize}
\tightlist
\item
  S1 Appendix.
\end{itemize}

\hypertarget{references}{%
\section*{References}\label{references}}
\addcontentsline{toc}{section}{References}

\hypertarget{refs}{}
\leavevmode\hypertarget{ref-dercleIdentificationNonSmall2020}{}%
1. Dercle L, Fronheiser M, Lu L, Du S, Hayes W, Leung DK, et al.
Identification of NonSmall Cell Lung Cancer Sensitive to Systemic Cancer
Therapies Using Radiomics. Clin Cancer Res. American Association for
Cancer Research; 2020;26: 2151--2162.
doi:\href{https://doi.org/10.1158/1078-0432.CCR-19-2942}{10.1158/1078-0432.CCR-19-2942}

\leavevmode\hypertarget{ref-coudrayClassificationMutationPrediction2018}{}%
2. Coudray N, Ocampo PS, Sakellaropoulos T, Narula N, Snuderl M, Fenyö
D, et al. Classification and mutation prediction from nonSmall cell lung
cancer histopathology images using deep learning. Nature Medicine.
Nature Publishing Group; 2018;24: 1559--1567.
doi:\href{https://doi.org/10.1038/s41591-018-0177-5}{10.1038/s41591-018-0177-5}

\leavevmode\hypertarget{ref-wayMachineLearningDetects2018}{}%
3. Way GP, Sanchez-Vega F, La K, Armenia J, Chatila WK, Luna A, et al.
Machine Learning Detects Pan-cancer Ras Pathway Activation in The Cancer
Genome Atlas. Cell Reports. Elsevier; 2018;23: 172--180.e3.
doi:\href{https://doi.org/10.1016/j.celrep.2018.03.046}{10.1016/j.celrep.2018.03.046}

\leavevmode\hypertarget{ref-loiPIK3CAMutationsAssociated2010}{}%
4. Loi S, Haibe-Kains B, Majjaj S, Lallemand F, Durbecq V, Larsimont D,
et al. PIK3CA mutations associated with gene signature of low mTORC1
signaling and better outcomes in estrogen receptorPositive breast
cancer. PNAS. National Academy of Sciences; 2010;107: 10208--10213.
doi:\href{https://doi.org/10.1073/pnas.0907011107}{10.1073/pnas.0907011107}

\leavevmode\hypertarget{ref-way_2018_machine_cellreports}{}%
5. Way GP, Sanchez-Vega F, La K, Armenia J, Chatila WK, Luna A, et al.
Machine Learning Detects Pan-cancer Ras Pathway Activation in The Cancer
Genome Atlas. Cell Reports. Elsevier; 2018;23: 172--180.e3.
doi:\href{https://doi.org/10.1016/j.celrep.2018.03.046}{10.1016/j.celrep.2018.03.046}

\leavevmode\hypertarget{ref-cantleyPhosphoinositide3KinasePathway2002}{}%
6. Cantley LC. The Phosphoinositide 3-Kinase Pathway. Science. American
Association for the Advancement of Science; 2002;296: 1655--1657.
doi:\href{https://doi.org/10.1126/science.296.5573.1655}{10.1126/science.296.5573.1655}

\leavevmode\hypertarget{ref-andre_2019_alpelisib_nengljmeda}{}%
7. André F, Ciruelos E, Rubovszky G, Campone M, Loibl S, Rugo HS, et al.
Alpelisib for PIK3CA-Mutated, Hormone ReceptorPositive Advanced Breast
Cancer. New England Journal of Medicine. Massachusetts Medical Society;
2019;380: 1929--1940.
doi:\href{https://doi.org/10.1056/NEJMoa1813904}{10.1056/NEJMoa1813904}

\leavevmode\hypertarget{ref-jacobsenCgdsrRbasedAPI2019}{}%
8. Jacobsen A, Luna A. Cgdsr: R-based API for accessing the MSKCC cancer
genomics data server (CGDS). 2019.

\leavevmode\hypertarget{ref-kuhnRecipesPreprocessingTools2020}{}%
9. Kuhn M, Wickham H. Recipes: Preprocessing tools to create design
matrices. 2020.

\leavevmode\hypertarget{ref-friedmanRegularizationPathsGeneralized2010}{}%
10. Friedman J, Hastie T, Tibshirani R. Regularization Paths for
Generalized Linear Models via Coordinate Descent. J Stat Softw. 2010;33:
1--22.

\leavevmode\hypertarget{ref-saitoPrecisionRecallPlotMore2015}{}%
11. Saito T, Rehmsmeier M. The Precision-Recall Plot Is More Informative
than the ROC Plot When Evaluating Binary Classifiers on Imbalanced
Datasets. PLOS ONE. Public Library of Science; 10: e0118432.
doi:\href{https://doi.org/10.1371/journal.pone.0118432}{10.1371/journal.pone.0118432}

\leavevmode\hypertarget{ref-kuhnTidymodelsEasilyInstall2020}{}%
12. Kuhn M, Wickham H. Tidymodels: Easily install and load the
'tidymodels' packages. 2020.

\leavevmode\hypertarget{ref-leroithInsulinlikeGrowthFactor2003}{}%
13. LeRoith D, Roberts CT. The insulin-like growth factor system and
cancer. Cancer Letters. 2003;195: 127--137.
doi:\href{https://doi.org/10.1016/S0304-3835(03)00159-9}{10.1016/S0304-3835(03)00159-9}

\leavevmode\hypertarget{ref-heInteractionInsulinReceptor1996}{}%
14. He W, Craparo A, Zhu Y, O'Neill TJ, Wang L-M, Pierce JH, et al.
Interaction of Insulin Receptor Substrate-2 (IRS-2) with the Insulin and
Insulin-like Growth Factor I Receptors EVIDENCE FOR TWO DISTINCT
PHOSPHOTYROSINE-DEPENDENT INTERACTION DOMAINS WITHIN IRS-2. J Biol Chem.
American Society for Biochemistry and Molecular Biology; 1996;271:
11641--11645.
doi:\href{https://doi.org/10.1074/jbc.271.20.11641}{10.1074/jbc.271.20.11641}

\leavevmode\hypertarget{ref-cizkovaGeneExpressionProfiling2010}{}%
15. Cizkova M, Cizeron-Clairac G, Vacher S, Susini A, Andrieu C,
Lidereau R, et al. Gene Expression Profiling Reveals New Aspects of
PIK3CA Mutation in ERalpha-Positive Breast Cancer: Major Implication of
the Wnt Signaling Pathway. PLOS ONE. Public Library of Science; 5:
e15647.
doi:\href{https://doi.org/10.1371/journal.pone.0015647}{10.1371/journal.pone.0015647}

\leavevmode\hypertarget{ref-muCNAPEMachineLearning2019}{}%
16. Mu Q, Wang J. CNAPE: A Machine Learning Method for Copy Number
Alteration Prediction from Gene Expression. IEEE/ACM Transactions on
Computational Biology and Bioinformatics. 2019; 1--1.
doi:\href{https://doi.org/10.1109/TCBB.2019.2944827}{10.1109/TCBB.2019.2944827}

\leavevmode\hypertarget{ref-gengHiddenMarkovModel2008}{}%
17. Geng H, Ali HH, Chan WC. A Hidden Markov Model Approach for
Prediction of Genomic Alterations from Gene Expression Profiling. In:
Măndoiu I, Sunderraman R, Zelikovsky A, editors. Bioinformatics Research
and Applications. Berlin, Heidelberg: Springer; 2008. pp. 414--425.
doi:\href{https://doi.org/10.1007/978-3-540-79450-9_38}{10.1007/978-3-540-79450-9\_38}

\leavevmode\hypertarget{ref-heGeneSignaturesAssociated2020}{}%
18. He X, Qin C, Zhao Y, Zou L, Zhao H, Cheng C. Gene signatures
associated with genomic aberrations predict prognosis in neuroblastoma.
Cancer Communications. 2020;40: 105--118.
doi:\href{https://doi.org/10.1002/cac2.12016}{10.1002/cac2.12016}

\nolinenumbers


\end{document}

